\fsection{Relación con el ciclo formativo}

El presente proyecto integra conocimientos y competencias adquiridos a lo largo del ciclo formativo, abarcando varios módulos de forma transversal:

\begin{itemize}
	\item Sistemas secuenciales programables y Sistemas programables avanzados: la adaptación del proyecto de TIA Portal requiere el manejo de bloques de función (FB), bloques de datos (DB) y la estructura de un programa S7-1500.

	\item Comunicaciones industriales: la configuración de la red Profinet, la asignación de direcciones IP estáticas y la comunicación entre dispositivos heterogéneos (PLC, IPC, cámara y robot).

	\item Informática industrial: la comprensión de la arquitectura Docker en el IPC, el frontend web de gestión y la integración del modelo de IA en el entorno industrial.

	\item Integración de sistemas de automatización industrial: la puesta en marcha del sistema completo, coordinando todos los elementos hardware y software para que operen de forma conjunta.

	\item Digitalización aplicada a los sectores productivos: la aplicación de inteligencia artificial y visión artificial en un entorno productivo real se enmarca directamente en este módulo.

	\item Documentación técnica: la elaboración de la presente memoria.
\end{itemize}


